\section[第13章\q 财务自由，关键不在赚多少]{第13章\\ 财务自由，关键不在赚多少}
\label{sec:ch13}

周一早上，地铁纽约之眼的门打开后，柔伊就导入人流中，流过灰砖走道，来到的巨大空间。

在走过白色意大利大理石地面时，柔伊想，眼窗就好比一处安放摄影镜头的地方。因为你得先安好取景点，才可以看到风景。

大厅询问处的台子上放着大束鲜花。今天选用的是白玫瑰和白色圣母百合。

{\kaishu 取景是拍出好照片的重要一步，你站在哪里，决定了你会看到什么，也就取决了你照片的画面。你了解我的意思吗？}

“我了解。”她一边走，一边低声对自己说。

之后，她又第一百次的想起了老妈的话，“其实我还有许多想做的事……”

她步入了西街知名的电视墙。今天电视墙面上的画面是太阳从美国西南部雄伟的山群后升起，撒下一片令人惊艳、紫橘色的曙光。

墙上出现了广告，柔伊站下脚步，仔细看着广告文本：

{\kaishu 你知道如何让梦想成真吗？}

{\kaishu 1次1美元，慢慢买下来。}

她走上手扶梯，上了2楼，站到外面开放空间，回头看西街，太阳正对着她，照在她所在杂志社的办公大楼。

她移动目光，今天这个位置可以看到大楼全貌。

早上柔伊和人资部门的大卫碰了面。他当场就为她设置了401帐户。大卫还给了她几个相关的建议，如何让退休帐户有效运作。当晚回到家之后，柔伊又上线重新开列了两个帐户，一个是摄影课程帐号，一个是旅游帐户。她也依照大卫传授的，设好了自动转户，只要收到薪水，就会有部分金额直接转入摄影课程帐户和旅游帐户，整个设置时间也不过5分钟。

她设置的转帐金额并不高，她想以后随时可以再调整。

在告别父亲，赶回纽约的前夜，柔伊想到亨利给她的建议绝不局限在财务上，还有更多。

现在，柔伊对自己的人生意义也更清楚了，她了解到什么对她来说，才是更重要的。她明白，自己就像年轻时的亨利，每天好几个小时的工作收入都花在不需要的事情上，这些事没有让她离自己的人生目标更近。

一旦认识到这一点，柔伊便意识到，其实不用花许多钱，也可以过生活。真的不需要赚太多，而是要清楚知道如何规划赚到的钱。

在回家的路上，她传了消息给洁西卡：

{\kaishu 洁西卡，真的很抱歉，我决定不换工作了，我在自由塔这里很开心。}

事实是，柔伊其实很喜欢目前工作，喜欢做杂志社的事，也喜欢杂志社共事的人。

她只需要做一些改变。

在和大卫谈过之后，她又去见了芭芭拉，告诉老板，她日后想会多一些时间出门旅游，就像年轻时的亨利一样，来一趟高调的休假研习。也就是说，1年当中，有1季的特刊她会无法参与。但是她会随身带的笔电，旅行的时候也可以工作。她谈的时候，其实很担心芭芭拉会不同意她的提议。

说完之后，芭芭拉沉默了一下，然后耸了一下肩膀，露出招牌表情：“可以的，但有个条件。”

“什么条件？”柔伊问。

“你要寄明信片给我。”

这天晚上，柔伊又做了一个梦。

她在缅因州的海岸边，缓缓的划着一条小船，吃着篮子里的小野莓。

“你看。”妈妈指着远处说。“白头海雕。”老爸补充说：“美国秃鹫。”柔伊用手遮着阳光，看到这些翱翔的鸟儿展翅上飞，飞过知名的高塔，穿过云霄，直到天空。

她醒来后，躺在床上。看着半暗的公寓，她犹疑着这里怎么那样安静。花了整整1分钟，她才想明白。

这个地方没有改变，今晚和过去一样的安静，只是她脑中的繁杂噪音消失了。那是一种对未来的担忧。就像冰箱嗡嗡嗡的声音，虽然吵，却因为习惯了而没有感觉。直到有一天，嗡嗡声戛然而止时，才会意识到之前的搅扰。

她在半暗半明的夜里，微笑了。

从某个角度来说，什么事都没有改变。不是设好了自动转帐后，帐户上就多出了许多存款来。但是她的感觉不一样了。因为知道，它们会成长的，枝上生枝，叶上生叶，积年累月，她的烦恼会自然消失了。

柔伊心情释然，转身又入睡了。一直到早上才醒。醒来之后，她从来没有感觉到这样的轻松，简直跟睡死过去似的。

不，应该修饰一下，是睡得轻松又自在。
